\documentclass[9pt]{article}
\bibliographystyle{alpha}
 
\usepackage{times}
\usepackage{transparency}
\usepackage[german]{babel}
\usepackage[T1]{fontenc}
\usepackage[latin1]{inputenc}


\screensize{8.5truein}{11truein}
\begin{document}

\title{LDNSD - A Fast and Scalable DNS Server}

\setbackground{background.png}
\begin{picture}(0,0)
	\put(315,-475){\includegraphics[scale=.3]{figures/protocollabs-font}}
\end{picture}
\subtitle{ }

\footerline{LDNSD - A Fast and Scalable DNS Server}
\date{{\tiny http://www.protocollabs.com}}
\name{Hagen Paul Pfeifer}
\email{\small{hagen.pfeifer@protcollabs.com}}
\company{ }
\maketitle

%%%%%%%%%%%%%%%%%%%%%%%%%%%%%%
\begin{slide}
\slidetitle{History DNS}{History DNS}
\bi
	\item In the early days (host)names to IP addresses are was achieved by manually maintained tables (host files)
	\item DNS was introduced to provide the same Information automatically through recursive queries
\ei
\end{slide}


%%%%%%%%%%%%%%%%%%%%%%%%%%%%%%
\begin{slide}
\slidetitle{Zones}{Zones}
\bi
	\item x
\ei
\end{slide}


%%%%%%%%%%%%%%%%%%%%%%%%%%%%%%
\begin{slide}
\slidetitle{Recursive - Iterative Nameserver}{Recursive - Iterative Nameserver}
\bi
	\item x
\ei
\end{slide}

%%%%%%%%%%%%%%%%%%%%%%%%%%%%%%
\begin{slide}
\slidetitle{Dynamic DNS}{Dynamic DNS}
\bi
	\item Mechanism to change DNS Information in real-time
	\item Standardized Mechanism: RFC 2136
	\item Angedacht um DHCP clients mit auffinbar zu machen
\ei
\end{slide}


%%%%%%%%%%%%%%%%%%%%%%%%%%%%%%
\begin{slide}
\slidetitle{Dynamic DNS}{Dynamic DNS}
\bi
	\item Wird bereits in vielen Gebieten eingesetzt - nur weiss man es oft nicht
	\item Standard Mechanismus (abfragen kann man DNS Informationen ohne zusaetzliche Software!)
	\item Keine propritaere Protokolle notwendig!
	\item DNS "sitzt" logisch unter IP. Sprich alle Dienste (egal ob WWW, FTP, XMPP, ...) funktionieren ohne Eingriff! (LDAP
	      liefert dies nicht, aber LDAP hat oft Funktionalitaet um DDNS Eintraege zu versenden!)
	\item DNS liefert keine Usermobilitaet, sondern eine Hostmobilitaet
\ei
\end{slide}


%%%%%%%%%%%%%%%%%%%%%%%%%%%%%%
\begin{slide}
\slidetitle{Dynamic DNS Clients}{Dynamic DNS Clients}
\bi
	\item nsupdate
	\item Python, Perl, Ruby modules
\begin{verbatim}
# nsupdate
> server ns.mydns.com
> update delete oldhost.example.com. A
> update add newhost.example.com. 86400 A 192.168.254.117
> send
> quit
\end{verbatim}
\ei
\end{slide}


%%%%%%%%%%%%%%%%%%%%%%%%%%%%%%
\begin{slide}
\slidetitle{Dynamic DNS Security}{Dynamic DNS Security}
\bi
	\item TSIG
\ei
\end{slide}

%%%%%%%%%%%%%%%%%%%%%%%%%%%%%%
\begin{slide}
\slidetitle{Multicast DNS}{Multicast DNS}
\bi
	\item Provide mechanism to operate in less configured environments (e.g. lookup a DNS a resource record in absence of a
	conventional managed DNS server)
\ei
\end{slide}


%%%%%%%%%%%%%%%%%%%%%%%%%%%%%%
\begin{slide}
\slidetitle{Nameserver Synchronization}{Nameserver Synchronization}
\bi
	\item Normally AXFR/IXFR Records
	\item TCP
\ei
\end{slide}



%%%%%%%%%%%%%%%%%%%%%%%%%%%%%%
\begin{slide}
\slidetitle{x}{x}
\bi
	\item x
\ei
\end{slide}



%%%%%%%%%%%%%%%%%%%%%%%%%%%%%%
\begin{slide}
\slidetitle{x}{x}
\bi
	\item x
\ei
\end{slide}


%%%%%%%%%%%%%%%%%%%%%%%%%%%%%%
\begin{slide}
\slidetitle{x}{x}
\bi
	\item x
\ei
\end{slide}

%%%%%%%%%%%%%%%%%%%%%%%%%%%%%%
\begin{slide}
\slidetitle{x}{x}
\bi
	\item x
\ei
\end{slide}

%%%%%%%%%%%%%%%%%%%%%%%%%%%%%%
\begin{slide}
\slidetitle{x}{x}
\bi
	\item x
\ei
\end{slide}

%%%%%%%%%%%%%%%%%%%%%%%%%%%%%%
\begin{slide}
\slidetitle{x}{x}
\bi
	\item x
\ei
\end{slide}

%%%%%%%%%%%%%%%%%%%%%%%%%%%%%%
\begin{slide}
\slidetitle{x}{x}
\bi
	\item x
\ei
\end{slide}


%%%%%%%%%%%%%%%%%%%%%%%%%%%%%%
\begin{slide}
\slidetitle{x}{x}
\bi
	\item x
\ei
\end{slide}

%%%%%%%%%%%%%%%%%%%%%%%%%%%%%%
\begin{slide}
\slidetitle{x}{x}
\bi
	\item x
\ei
\end{slide}


%%%%%%%%%%%%%%%%%%%%%%%%%%%%%%
\begin{slide}
\slidetitle{x}{x}
\bi
	\item x
\ei
\end{slide}


%%%%%%%%%%%%%%%%%%%%%%%%%%%%%%
\begin{slide}
\slidetitle{x}{x}
\bi
	\item x
\ei
\end{slide}




\end{document}

% vim600: fdm=marker tw=130 sw=4 ts=4 sts=4 ff=unix noet:
